\begin{table}[!ht]
\centering
\begin{tabular}{r l l}
\toprule
Byte offset & Encoded bytes & Instruction\\
\midrule
0 & \texttt{01 00} & PUSH8 0\\
2 & \texttt{01 03} & PUSH8 3\\
4 & \texttt{20} & DUP\\
5 & \texttt{31 10} & JZ 16\\
7 & \texttt{22} & SWAP\\
8 & \texttt{23} & OVER\\
9 & \texttt{10} & ADD\\
10 & \texttt{22} & SWAP\\
11 & \texttt{01 01} & PUSH8 1\\
13 & \texttt{11} & SUB\\
14 & \texttt{30 04} & JMP 4\\
16 & \texttt{21} & DROP\\
17 & \texttt{00} & HALT\\
\bottomrule
\end{tabular}
\caption{The 18-byte countdown image sums three, two and one. The instruction at offset five has target byte hexadecimal \texttt{10}, decimal 16. Offset six is its operand, even though its byte value is also the ADD opcode. Instruction identity comes from decoding at an established start.}
\label{tab:revision-bytecode-countdown}
\end{table}
